% !TEX program = xelatex
\documentclass[11pt,a4paper,landscape]{article}

\usepackage[UTF8,heading=true]{ctex}
\setCJKmainfont{SimSun}[BoldFont=SimHei,ItalicFont=KaiTi]
\setCJKsansfont{SimHei}
\setCJKmonofont{FangSong}
\usepackage[top=12mm,bottom=12mm,left=12mm,right=12mm]{geometry}
\pagestyle{empty}

\usepackage[dvipsnames,svgnames,x11names]{xcolor}
\definecolor{primary}{HTML}{1E5C44}
\definecolor{primaryLight}{HTML}{5AA07A}
\definecolor{primaryDark}{HTML}{123C2C}
\definecolor{accent}{HTML}{D97904}
\definecolor{accentLight}{HTML}{F2B84B}
\definecolor{danger}{HTML}{B63A2B}
\definecolor{purple}{HTML}{54428E}
\definecolor{teal}{HTML}{0B7285}
\definecolor{bgPage}{HTML}{F4F1E8}
\definecolor{bgCard}{HTML}{FFFDF8}
\definecolor{gray80}{HTML}{3A3A3A}
\definecolor{gray60}{HTML}{6B6B6B}
\definecolor{gray20}{HTML}{DDD7C8}

\usepackage{tikz}
\usetikzlibrary{arrows.meta,positioning,shadows.blur,calc,shapes.geometric,fit,backgrounds}
\usepackage{fontawesome5}
\usepackage{graphicx}
\usepackage{setspace}
\usepackage{anyfontsize}
\usepackage{hyperref}
\hypersetup{colorlinks=false}

\tikzset{
  card/.style={
    rectangle, rounded corners=6pt, draw=gray20, line width=0.4pt, fill=bgCard,
    inner sep=8pt, align=center,
    blur shadow={shadow blur steps=6,shadow xshift=0.6pt,shadow yshift=-0.8pt,
      shadow blur radius=2pt,shadow opacity=12}
  },
  bigIconCircle/.style={
    circle, minimum size=18mm, fill=#1, text=white,
    font=\LARGE, inner sep=0pt
  },
  phase/.style={
    rectangle, rounded corners=4pt, minimum width=145mm, minimum height=16mm,
    fill=#1, text=white, font=\sffamily\bfseries\normalsize,
    inner xsep=10pt, align=left
  }
}

\newcommand{\pagemark}[1]{%
  \begin{tikzpicture}[remember picture,overlay]
    \node[anchor=south east,font=\sffamily\tiny\color{gray60}]
      at ([xshift=-8mm,yshift=8mm]current page.south east) {#1};
  \end{tikzpicture}%
}

\begin{document}

% Page 1 Cover
\begin{tikzpicture}[remember picture,overlay]
  \shade[left color=primaryDark,right color=primaryLight]
    (current page.south west) rectangle (current page.north east);
  \fill[white,opacity=0.05] ([xshift=6cm,yshift=-2.5cm]current page.north east) circle(8cm);
  \fill[accent,opacity=0.14] ([xshift=-9cm,yshift=1.5cm]current page.south west) circle(5cm);
  \node[text width=21cm,align=center] at ([yshift=2.2cm]current page.center) {
    {\color{white}\sffamily\fontsize{34}{40}\selectfont\bfseries 从“中国向何处去”}\\[4pt]
    {\color{white}\sffamily\fontsize{34}{40}\selectfont\bfseries 到“建设一个新中国”}\\[10pt]
    {\color{white}\sffamily\fontsize{20}{26}\selectfont 《新民主主义论》读书报告}
  };
  \draw[white,opacity=0.4,line width=0.8pt]
    ([xshift=-7cm,yshift=-0.3cm]current page.center)--([xshift=7cm,yshift=-0.3cm]current page.center);
  \node at ([yshift=-1.7cm]current page.center) {
    {\color{white}\sffamily\small \faCalendar\; 2026年4月18日 \qquad \faUser\; 姓名/学号待填写}
  };

  \begin{scope}[shift={([xshift=-6cm,yshift=-4.6cm]current page.center)}]
    \fill[white,opacity=0.14,rounded corners=8pt] (-31mm,-13mm) rectangle (31mm,13mm);
    \node[align=center] at (0,0) {
      {\color{white}\sffamily\bfseries\fontsize{26}{30}\selectfont 1 个核心问题}\\[2pt]
      {\color{white}\small 中国向何处去}
    };
  \end{scope}
  \begin{scope}[shift={([xshift=0cm,yshift=-4.6cm]current page.center)}]
    \fill[white,opacity=0.14,rounded corners=8pt] (-31mm,-13mm) rectangle (31mm,13mm);
    \node[align=center] at (0,0) {
      {\color{white}\sffamily\bfseries\fontsize{26}{30}\selectfont 3 条主线}\\[2pt]
      {\color{white}\small 国情、道路、重建}
    };
  \end{scope}
  \begin{scope}[shift={([xshift=6cm,yshift=-4.6cm]current page.center)}]
    \fill[white,opacity=0.14,rounded corners=8pt] (-31mm,-13mm) rectangle (31mm,13mm);
    \node[align=center] at (0,0) {
      {\color{white}\sffamily\bfseries\fontsize{26}{30}\selectfont 1 个结论}\\[2pt]
      {\color{white}\small 新中国不是旧路延长线}
    };
  \end{scope}
\end{tikzpicture}
\pagemark{1 / 6}
\clearpage

% Page 2
\begin{tikzpicture}[remember picture,overlay]
  \fill[primary] (current page.north west) rectangle ([yshift=-2.2cm]current page.north east);
  \node[anchor=west,font=\sffamily\LARGE\bfseries\color{white}]
    at ([xshift=2cm,yshift=-1.1cm]current page.north west) {\faCompass\; 全文最先回答什么问题};
  \fill[bgPage] ([yshift=-2.2cm]current page.north west) rectangle (current page.south east);

  \node[card,minimum width=10.5cm,minimum height=10cm] (q) at ([xshift=-6.4cm,yshift=0.3cm]current page.center) {};
  \node[bigIconCircle=accent] at ([yshift=2.7cm]q.center) {\faQuestion};
  \node[font=\sffamily\bfseries\fontsize{24}{28}\selectfont\color{gray80}] at ([yshift=0.8cm]q.center) {中国向何处去};
  \node[font=\sffamily\large\color{gray60},text width=8.7cm,align=left] at ([yshift=-1.8cm]q.center) {
    这不是书斋里的抽象追问，而是1940年前后关于民族存亡、革命道路与国家前途的现实问题。
  };

  \node[card,minimum width=13cm,minimum height=10cm] (ans) at ([xshift=6.2cm,yshift=0.3cm]current page.center) {};
  \node[font=\sffamily\bfseries\Large\color{primaryDark}] at ([yshift=3.2cm]ans.center) {毛泽东的回答方式};
  \node[font=\sffamily\large\color{gray80},text width=11cm,align=left] at ([yshift=0.6cm]ans.center) {
    1. 先判断中国社会的性质\\[7pt]
    2. 再分析革命所处的历史阶段\\[7pt]
    3. 最后推出中国应走的新民主主义道路
  };
  \node[rectangle,rounded corners=6pt,fill=accent!12,draw=accent!50,minimum width=11cm,minimum height=2.2cm,align=center]
    at ([yshift=-3.2cm]ans.center) {{\sffamily\bfseries\large\color{accent} 先认识国情，再决定道路}};
\end{tikzpicture}
\pagemark{2 / 6}
\clearpage

% Page 3
\begin{tikzpicture}[remember picture,overlay]
  \fill[accent] (current page.north west) rectangle ([yshift=-2.2cm]current page.north east);
  \node[anchor=west,font=\sffamily\LARGE\bfseries\color{white}]
    at ([xshift=2cm,yshift=-1.1cm]current page.north west) {\faStar\; “新”这个字到底新在哪里};
  \fill[bgPage] ([yshift=-2.2cm]current page.north west) rectangle (current page.south east);

  \node[card,minimum width=7cm,minimum height=9cm] (c1) at ([xshift=-8.2cm,yshift=0cm]current page.center) {};
  \node[bigIconCircle=primary] at ([yshift=2.2cm]c1.center) {\faClock};
  \node[font=\sffamily\bfseries\Large\color{gray80}] at ([yshift=0.5cm]c1.center) {时代变了};
  \node[font=\sffamily\small\color{gray60},text width=5.8cm,align=left] at ([yshift=-2.1cm]c1.center) {它发生在帝国主义时代，不再属于旧世界资产阶级民主革命的逻辑。};

  \node[card,minimum width=7cm,minimum height=9cm] (c2) at ([xshift=0cm,yshift=0cm]current page.center) {};
  \node[bigIconCircle=teal] at ([yshift=2.2cm]c2.center) {\faUsers};
  \node[font=\sffamily\bfseries\Large\color{gray80}] at ([yshift=0.5cm]c2.center) {领导变了};
  \node[font=\sffamily\small\color{gray60},text width=5.8cm,align=left] at ([yshift=-2.1cm]c2.center) {能把民族独立与社会变革真正结合起来的，是无产阶级及其政党。};

  \node[card,minimum width=7cm,minimum height=9cm] (c3) at ([xshift=8.2cm,yshift=0cm]current page.center) {};
  \node[bigIconCircle=purple] at ([yshift=2.2cm]c3.center) {\faRoad};
  \node[font=\sffamily\bfseries\Large\color{gray80}] at ([yshift=0.5cm]c3.center) {前途变了};
  \node[font=\sffamily\small\color{gray60},text width=5.8cm,align=left] at ([yshift=-2.1cm]c3.center) {它不是停在资本主义，而是把社会主义作为发展方向。};

  \node[rectangle,rounded corners=6pt,fill=primary!10,draw=primary!40,line width=0.6pt,minimum width=24cm,minimum height=1.8cm,align=center]
    at ([yshift=-5.2cm]current page.center) {{\sffamily\bfseries\large\color{primaryDark} 所谓“新”，不是换个词，而是历史条件改变后的重新定位。}};
\end{tikzpicture}
\pagemark{3 / 6}
\clearpage

% Page 4
\begin{tikzpicture}[remember picture,overlay]
  \fill[teal] (current page.north west) rectangle ([yshift=-2.2cm]current page.north east);
  \node[anchor=west,font=\sffamily\LARGE\bfseries\color{white}]
    at ([xshift=2cm,yshift=-1.1cm]current page.north west) {\faLayerGroup\; 新中国不是只换政权，而是整体重建};
  \fill[bgPage] ([yshift=-2.2cm]current page.north west) rectangle (current page.south east);

  \node[phase=primary,anchor=north] at ([yshift=-3.1cm]current page.north) {政治、经济、文化不是三个散点，而是一整套社会方案};

  \node[card,minimum width=7.2cm,minimum height=7.4cm] (p) at ([xshift=-8.2cm,yshift=-0.6cm]current page.center) {};
  \node[bigIconCircle=primaryDark] at ([yshift=1.8cm]p.center) {\faLandmark};
  \node[font=\sffamily\bfseries\Large\color{gray80}] at ([yshift=0.2cm]p.center) {新政治};
  \node[font=\sffamily\small\color{gray60},text width=6cm,align=left] at ([yshift=-2.0cm]p.center) {核心是新民主主义政治，回答的是“谁来领导、如何建立新国家”。};

  \node[card,minimum width=7.2cm,minimum height=7.4cm] (e) at ([xshift=0cm,yshift=-0.6cm]current page.center) {};
  \node[bigIconCircle=accent] at ([yshift=1.8cm]e.center) {\faIndustry};
  \node[font=\sffamily\bfseries\Large\color{gray80}] at ([yshift=0.2cm]e.center) {新经济};
  \node[font=\sffamily\small\color{gray60},text width=6cm,align=left] at ([yshift=-2.0cm]e.center) {核心是新民主主义经济，回答的是“社会基础怎样重建”。};

  \node[card,minimum width=7.2cm,minimum height=7.4cm] (c) at ([xshift=8.2cm,yshift=-0.6cm]current page.center) {};
  \node[bigIconCircle=purple] at ([yshift=1.8cm]c.center) {\faBook};
  \node[font=\sffamily\bfseries\Large\color{gray80}] at ([yshift=0.2cm]c.center) {新文化};
  \node[font=\sffamily\small\color{gray60},text width=6cm,align=left] at ([yshift=-2.0cm]c.center) {“民族的、科学的、大众的文化”回答的是“培养什么样的人与社会精神”。};
\end{tikzpicture}
\pagemark{4 / 6}
\clearpage

% Page 5
\begin{tikzpicture}[remember picture,overlay]
  \fill[purple] (current page.north west) rectangle ([yshift=-2.2cm]current page.north east);
  \node[anchor=west,font=\sffamily\LARGE\bfseries\color{white}]
    at ([xshift=2cm,yshift=-1.1cm]current page.north west) {\faBookOpen\; 和《中国革命和中国共产党》一起读};
  \fill[bgPage] ([yshift=-2.2cm]current page.north west) rectangle (current page.south east);

  \node[rectangle,rounded corners=6pt,fill=danger!8,draw=danger!70,minimum width=10.5cm,minimum height=1.6cm]
    (leftHdr) at ([xshift=-6.3cm,yshift=3.7cm]current page.center)
    {{\sffamily\bfseries\large\color{danger} \faMap\; 《新民主主义论》}};
  \node[rectangle,rounded corners=6pt,fill=primary!10,draw=primary!70,minimum width=10.5cm,minimum height=1.6cm]
    (rightHdr) at ([xshift=6.3cm,yshift=3.7cm]current page.center)
    {{\sffamily\bfseries\large\color{primaryDark} \faProjectDiagram\; 《中国革命和中国共产党》}};

  \node[card,minimum width=10.5cm,minimum height=8.2cm] at ([xshift=-6.3cm,yshift=-0.1cm]current page.center) {
    \begin{minipage}{9.1cm}
    \centering
    {\sffamily\bfseries\large\color{gray80} 更像目标与方案}\\[8pt]
    \raggedright
    {\sffamily\small\color{gray60}
    重点回答：\\[5pt]
    • 新中国应是什么样\\
    • 为什么要走新民主主义道路\\
    • 新政治、新经济、新文化如何构成整体}
    \end{minipage}
  };

  \node[card,minimum width=10.5cm,minimum height=8.2cm] at ([xshift=6.3cm,yshift=-0.1cm]current page.center) {
    \begin{minipage}{9.1cm}
    \centering
    {\sffamily\bfseries\large\color{gray80} 更像现实依据与分析}\\[8pt]
    \raggedright
    {\sffamily\small\color{gray60}
    重点回答：\\[5pt]
    • 中国社会性质是什么\\
    • 革命对象、任务、动力是什么\\
    • 现阶段革命的性质与前途如何判断}
    \end{minipage}
  };

  \node[rectangle,rounded corners=6pt,fill=accent!12,draw=accent!50,line width=0.6pt,minimum width=24.5cm,minimum height=1.9cm,align=center]
    at ([yshift=-5.2cm]current page.center) {{\sffamily\bfseries\large\color{accent} 一篇像路线图，一篇像地形图；两篇一起读，才更能看出“从国情推出道路”的方法。}};
\end{tikzpicture}
\pagemark{5 / 6}
\clearpage

% Page 6
\begin{tikzpicture}[remember picture,overlay]
  \shade[left color=accent,right color=primary]
    (current page.south west) rectangle (current page.north east);
  \fill[white,opacity=0.05] ([xshift=-6cm,yshift=3cm]current page.south west) circle(7.5cm);
  \fill[white,opacity=0.03] ([xshift=7cm,yshift=-3cm]current page.north east) circle(9cm);

  \node[text width=21cm,align=center] at ([yshift=1.8cm]current page.center) {
    {\color{white}\sffamily\fontsize{30}{36}\selectfont\bfseries 我读完后的最大收获}\\[14pt]
    {\color{white}\sffamily\fontsize{20}{28}\selectfont 理论最有力量的时候，}\\[4pt]
    {\color{white}\sffamily\fontsize{20}{28}\selectfont 不是概念最复杂的时候，}\\[4pt]
    {\color{white}\sffamily\fontsize{20}{28}\selectfont 而是它最能回应现实问题的时候。}
  };

  \node[rectangle,rounded corners=8pt,fill=white,opacity=0.14,minimum width=22cm,minimum height=2.4cm] at ([yshift=-3.6cm]current page.center) {};
  \node[text width=20cm,align=center] at ([yshift=-3.6cm]current page.center) {
    {\color{white}\sffamily\large “新”不是口号式的新，而是在历史条件、领导力量和发展方向都变化之后，对中国革命道路作出的创造性回答。}
  };
\end{tikzpicture}
\pagemark{6 / 6}
\clearpage

\end{document}
